\begin{lemma}[Change lemma]
Fix a finite loopless multigraph with its exact edge set, a palette of
$k$ colors, and a distinguished edge $e_0$.  Let $(T,\varphi_n)$ be a
positive-rung extended Tashkinov tree.  Let $\mathcal B$ be a hierarchy
base and let
\[
 T_n=T_{n,0}\subset T_{n,1}\subset\cdots\subset
 T_{n,q}\subset T_{n,q+1}=T
\]
be a good hierarchy under $\varphi_n$, where $q$ is its last divider
index and
\[
 A=T_{n,q}^{*}=\noderef{HierarchyStarredStage}(H,q).
\]
Assume that $T$ has maximum property, that $V(T)$ is not elementary,
that $\varphi_n$ is proper and leaves exactly $e_0$ uncolored, and that
the rung number is positive.  These data form a
\noderef{HierarchyCounterexample} $C$.

Assume the six lower-rung assertions of the technical theorem, uniformly
for every positive-rung maximum ETT on the same graph, palette, and
distinguished edge, with arbitrary certified input and generating
colorings and with rung number smaller than that of $T$: lower-rung
elementarity, terminal interchangeability, the two PE stability and
stored-pair assertions, the terminal PE component assertion, the
defective-endpoint order assertion, and the mod-coloring assertion.
Assume also the same-rung induction hypothesis that every maximum ETT
with a good hierarchy having fewer than $q$ dividers is elementary.
Finally assume the large-color estimate
\[
 2n+11\leq |\overline\varphi_n(T_n)|.
\]

Put $p=\noderef{HierarchyPathNumber}(H)$ and let
\[
 T(y_{p-1})=T\big|_{\min\{|A|+p-1,|T|\}}
\]
be the exact ordered prefix supplied by \noderef{TreeSequencePrefix}.
Assume (6.6): $V(T(y_{p-1}))$ is elementary under $\varphi_n$.
Assume moreover that $C$ is lexicographically minimal among all
\noderef{HierarchyCounterexample}s on the fixed graph, palette, and
$e_0$ which have exactly $q$ dividers, allowing their certified input
colorings, generating colorings, bases, ETTs, and good hierarchies to
vary.  Explicitly, no such counterexample $C'$ satisfies
\[
 \noderef{HierarchyCounterexampleLess}
  \bigl(\noderef{HierarchyCounterexampleMeasure}(C'),
        \noderef{HierarchyCounterexampleMeasure}(C)\bigr).
\]

Let $\alpha\ne\beta$ be colors missing on $T(y_{p-1})$.  Write
\[
 i_c=\noderef{HierarchyFirstMissingPosition}(\varphi_n,T,c),
 \qquad v_c=T.\mathit{vertex}(i_c).
\]
If at least one of $\alpha,\beta$ is not missing on $A$, assume both
\[
 v_\alpha\prec v_\beta                                      \tag{a}
\]
and the exact edge-identity exclusion
\[
 \alpha\notin
 \left\{c:\exists t\in\operatorname{Fin}|T|,
       |A|\leq t<i_\beta,
       \ \varphi_n(T.\mathit{edge}(t))=c\right\}.             \tag{b}
\]
Thus (b) is precisely
$\alpha\notin\varphi_n\langle T(v_\beta)-A\rangle$; it is not
an endpoint-only abbreviation.

Then the full two-color components through $v_\alpha$ and $v_\beta$
are equal in each of the following cases.
\begin{enumerate}
\item $q\geq1$, and either
$\alpha\in\overline\varphi_n(T_{n,q})$ or
$\{\alpha,\beta\}\cap D_{n,q}=\varnothing$.
\item $q=0$, and either
$\alpha\in\overline\varphi_n(T_n)$ or
$\{\alpha,\beta\}\cap D_n=\varnothing$.
\item $\alpha\in\overline\varphi_n(A)$ and the singleton
$\{\alpha\}$ is closed on $A$ under $\varphi_n$.
\end{enumerate}
In case 3 the colors are also $A$-interchangeable.  Consequently their
common nontrivial $(\alpha,\beta)$-path component is the only such path
component meeting $A$.
\end{lemma}
\begin{proof}
Unpack $C$.  Its fields give positivity, the displayed
\noderef{TashkinovHierarchy}, its \noderef{GoodHierarchy} witness,
\noderef{MaximumProperty}, and non-elementarity.  Let
$K_x$ denote the full \noderef{TwoColorComponent} through the vertex
node $x$.  All components below therefore contain the actual edge nodes
of the multigraph; parallel edges are never identified by their endpoint
pairs.

We first record consequences of (6.6).  The last hierarchy stage is a
proper ordered initial segment of $T$.  By the definition
\noderef{HierarchyPathNumber}, $p\geq1$.  Hence $A$, every starred
stage, and every ordinary hierarchy stage are ordered initial segments
of $T(y_{p-1})$.  This follows entry by entry from
\noderef{TashkinovHierarchy},
\noderef{HierarchyStarredStage}, and
\noderef{TreeSequenceInitialSegment}.  Restricting the pairwise
disjoint missing-color condition in
\noderef{ColoredVertexSetElementary} therefore shows that all these
stages, including $A$, are elementary under $\varphi_n$.               \tag{1}

Because $\alpha$ and $\beta$ are missing on this elementary prefix,
the least-candidate clause of
\noderef{HierarchyFirstMissingPosition} makes $v_\alpha$ and
$v_\beta$ actual vertices of the prefix missing the indicated colors.
They are distinct: otherwise the same vertex would miss the two colors,
but $\alpha\ne\beta$ and this does not yet contradict elementarity;
instead, if they coincided then $K_{v_\alpha}=K_{v_\beta}$ immediately,
so every asserted component equality is already proved.  We may thus
assume that they are distinct.  Elementarity says that $\beta$, which is
missing at $v_\beta$, is not missing at $v_\alpha$, and symmetrically
$\alpha$ is not missing at $v_\beta$.  By the complementary definitions
\noderef{MissingColors} and \noderef{PresentColors}, each vertex is
incident with an edge of the other color.  Properness makes it a
degree-one end in the two-color incidence graph.  Hence $K_{v_\alpha}$
and $K_{v_\beta}$ are nontrivial path components in the sense of
\noderef{TwoColorPathComponent}.  This verifies the path-component
premise each time \noderef{TreeInterchangeable} is used below.           \tag{2}

We shall repeatedly use the following endpoint observation.  Under the
proper coloring $\varphi_n$, every two-color component has a simple
spanning path by \noderef{TwoColorComponentSpanningPath}; the incidence
degree is at most two by
\noderef{TwoColorIncidenceDegreeAtMostTwo}.  An end vertex of a
nontrivial such path is a graph vertex missing one of its two colors:
an edge-node has two incidence neighbors because the graph is loopless,
whereas at a graph vertex properness permits at most one edge of either
color, and degree one means that the other color is absent.  Thus a
nontrivial component which cannot cross the boundary of an elementary
set has both ends among the unique vertices of that set missing its two
colors.  Distinct components are vertex-disjoint by
\noderef{TwoColorComponentsDisjoint}.                                  \tag{3}

We next make explicit the preservation argument used whenever a
component through $v_\beta$ is changed.  Suppose $K_{v_\alpha}\ne
K_{v_\beta}$ and set
\[
 \sigma_n=\varphi_n/K_{v_\beta}.
\]
By \noderef{TwoColorComponentSwap} the swap changes only the exact
edge nodes of $K_{v_\beta}$ and only interchanges $\alpha$ and $\beta$.
It fixes the value $\mathsf{none}$, and
\noderef{TwoColorComponentSwapProper} makes $\sigma_n$ proper.
Consequently \noderef{IsOneEdgeUncolored} shows that precisely $e_0$
remains uncolored.  Since $v_\beta$ misses $\beta$ before the swap, it
misses $\alpha$ afterwards; the different component through
$v_\alpha$ is untouched and $v_\alpha$ still misses $\alpha$.
Therefore any certified prefix containing both vertices is
non-elementary under $\sigma_n$.                                       \tag{4}

Here is the exact smaller-counterexample test which will be cited in the
three location cases.  Let $B$ be a hierarchy stage or starred stage
which is an ordered prefix of the prefix $T(v_\beta)$, and suppose:
\begin{enumerate}
\item the swap is stable on $B$ with every control set required by its
hierarchy presentation, and every displayed edge identity of $B$ keeps
its value;
\item every displayed edge identity of $T(v_\beta)-B$ has the same
value under $\sigma_n$ and $\varphi_n$;
\item every vertex of $T(v_\beta)-v_\beta$ has the same missing-color
set under the two colorings;
\item the corresponding ETT, maximum property, hierarchy base, and all
five good-hierarchy fields have been verified under $\sigma_n$.
\end{enumerate}
Then these data and (4) form an actual
\noderef{HierarchyCounterexample} $C'$ with the same $q$ dividers.
The terminal stage is the exact prefix $T(v_\beta)$, not an abstract
tree having the same vertex set.

We also check its measure.  The vertex $v_\beta$ belongs to
$T(y_{p-1})$.  If it lies outside $A$, then its index is at most
$|A|+p-1$.  The terminal tail of $T(v_\beta)$ after $A$ consequently
has length at most $p-1$.  In
\noderef{HierarchyPathNumber}, the fallback is that tail length and
every candidate lies between $1$ and that length.  Thus the new path
number is at most $p-1<p$, which strictly lowers the first coordinate of
\noderef{HierarchyCounterexampleMeasure}.  Therefore
\noderef{HierarchyCounterexampleLess} gives
\noderef{HierarchyCounterexampleLess} give
\[
                  \mathfrak m(C')<_{\rm hier}\mathfrak m(C). \tag{5}
\]
The assumed minimality forbids $C'$.  Notice that (5) is applied only
after the actual base, hierarchy, maximum property, non-elementarity,
path number, and terminal-tail coordinate have all been checked.

We now distinguish the locations of
$a=v_\alpha$ and $b=v_\beta$.  If at least one is outside $A$, the
conditional hypotheses (a)--(b) apply.  Since $A$ is an ordered initial
segment, (a) rules out $a\notin A$, $b\in A$.  The remaining
possibilities are exactly the following three cases.

\textbf{Location 1: $a,b\in V(A)$.}
By (1), $A$ is elementary; hence $a$ and $b$ are respectively the
unique vertices of $A$ missing $\alpha$ and $\beta$.  If both singleton
colors are closed on $A$, no exact boundary edge identity of $A$ is
colored $\alpha$ or $\beta$, by
\noderef{TreeColorSetClosed} and
\noderef{BoundaryColorEdges}.  The component through $a$ cannot leave
$A$.  Observation (3) says that its other path end must be $b$, so
$K_a=K_b$.  The same argument applied to an arbitrary nontrivial
$(\alpha,\beta)$-path component meeting $A$ forces its two ends to be
$a,b$; hence this common component is the only one meeting $A$.          \tag{6}

Suppose precisely one of the two colors is closed on $A$.  If $q>0$,
then $A=T_{n,q}$ and the last-divider instance of
\noderef{step1}, with the closed color placed in its first color slot,
applies.  Its lower-rung premise is exactly the six-assertion hypothesis
of this lemma, its fewer-divider premise is our same-rung induction
hypothesis, and its numerical premise is the displayed $2n+11$ estimate.
It makes the two colors $A$-interchangeable.  If $q=0$ and the terminal
kind is PE, use the core conclusion (iii) of
\noderef{interchange}; reflexivity of the two clauses in
\noderef{CombinedStableColoring} supplies its stability premise, and
our six lower-rung package supplies its induction premise.  If $q=0$
and the terminal kind is not PE, \noderef{HierarchyBase} identifies
$A$ with $T_n$.  Inspecting the SE or RE constructor in
\noderef{TashkinovSeriesStep}, and using symmetry of stability from
\noderef{sc2} in the SE case, shows that every color missing on $T_n$
is closed there.  Thus the assumed ``precisely one'' subcase cannot
occur.  Interchangeability in the two possible cases gives both
$K_a=K_b$ and uniqueness of the path component meeting $A$ by the
definition \noderef{TreeInterchangeable}.                              \tag{7}

It remains in this location to consider the case in which neither color
is closed on $A$.  Let $r$ be the first starred stage on which $\beta$
is missing.  Both $a$ and $b$ lie in that stage.  If $r>0$, the witness
is new above the preceding starred stage; (1) and
\noderef{GoodHierarchyNewMissingClosed} make $\beta$ closed on
$T_{n,r}=T_{n,r}^{*}$.  Repeating (7) with $r$ in place of $q$ gives
$K_a=K_b$.  If $r=0$ and the last extension is non-PE, the same SE/RE
closedness argument makes both colors closed on $T_n=A$, and (6)
applies.  If $r=0$ and the extension is PE, place $a$ and $b$ in the
two constituents $T_n$ and $R_n$ of the exact union $A$.  The four
conclusions (i), (ii), and (iv) of \noderef{interchange}, as appropriate
to those two locations, show that their two path components coincide.
This is the paper's descent to the first starred stage; it uses the PE
base interface rather than pretending that the union itself is a ladder
ETT.                                                                    \tag{8}

Equations (6)--(8) prove the required component equality in Location 1
under each of alternatives 1--3 in the statement.  They also prove the
additional $A$-interchangeability whenever alternative 3 holds.

\textbf{Location 2: $a,b\notin V(A)$.}
Now (a)--(b) hold.  Since neither color is missing on $A$, alternative 1
forces $\{\alpha,\beta\}\cap D_{n,q}=\varnothing$ when $q>0$, and
alternative 2 forces $\{\alpha,\beta\}\cap D_n=\varnothing$ when
$q=0$.  We first control the exact displayed identities of $A$.  By
\noderef{HkuTreeColorsControlled}, a color used on $T_n$ belongs to
$\overline\varphi_n(T_n)\cup D_n$.  Neither $\alpha$ nor $\beta$ is
missing on $T_n$, since $T_n$ is an initial segment of $A$.  Moreover
$D_{n,q}=D_n\setminus\overline\varphi_n(A)$ by
\noderef{HierarchyStoredColorsAt}; because neither color is missing on
$A$, membership in $D_n$ would be membership in $D_{n,q}$, which has
just been excluded.  Thus neither color is used on $T_n$.

Every displayed edge of $A-T_n$ and then of $T(b)-A$ is introduced by
the certified augmentation sequence.  By
\noderef{TashkinovAugmentation}, its exact color is missing on its old
prefix.  Such a color is missing on the later tree $A$, so it is not
$\alpha$ or $\beta$ in the first interval.  In the second interval, the
first-missing property of $b$ excludes $\beta$, while (b) excludes
$\alpha$.  The possible SE/RE connecting edge is controlled by the same
stored-color calculation.  Therefore
\[
 \alpha,\beta\notin
 \varphi_n\langle T(b)-A\rangle.                              \tag{9}
\]

Assume for contradiction that $K_a\ne K_b$ and perform the swap defining
$\sigma_n$.  A vertex of $A$ on $K_b$ misses neither $\alpha$ nor
$\beta$.  Properness gives exactly one incident edge of each color there,
and both edge nodes belong to the same full component.  The swap merely
exchanges these two incident values, so the vertex missing set is
unchanged.  If a changed edge is incident with $A$, its old color is
$\alpha$ or $\beta$; these colors are neither missing on $A$ nor members
of the relevant control set.  Thus the edge clause of
\noderef{StableColoring} imposes no equality on that edge, while every
controlled incident edge has another color and is fixed.  The swap is
therefore stable on $A$ even though $K_b$ need not be disjoint from
$A$.  The preceding calculation on $A$ together with (9) fixes every
displayed identity of $T(b)$.  At a vertex strictly before $b$, membership in the full
component again exchanges a present pair, while nonmembership fixes all
incident values; its missing set is unchanged.  At $b$ only $\beta$ is
replaced by $\alpha$.  This proves the edge and vertex assertions
(2)--(3) of the smaller-counterexample test.                              \tag{10}

We verify the certified data, rather than infer it from the unchanged
vertex set.  Restrict the final Tashkinov augmentation sequence in
\noderef{TashkinovSeriesStep} and
\noderef{TashkinovClosure} at the exact length of $T(b)$.  At each
retained step, \noderef{TashkinovAugmentation} names the new edge
identity, its color, and an old missing-color witness.  Equations
(9)--(10) preserve that same edge value and witness, so the same finite
sequence is valid under $\sigma_n$.  Restricting the terminal stable
relation to $T_n$ and applying \noderef{Td} gives a
\noderef{ModColoring}.  Its universal clause supplies an actual
\noderef{CorrespondingExtendedTashkinovTree} on $T(b)$, and
\noderef{MaximumPropertyTransport} transfers maximum property.          \tag{11}

Use the same closure and core in the PE case, and the terminal ladder
tree as core in the non-PE case, to construct the new
\noderef{HierarchyBase}.  Copy all ordinary stages through $T_{n,q}$,
the divider positions, and take $T(b)$ as the new terminal stage.
Equation (a) and $b\notin A$ make the last containment strict; all
earlier strict and initial-segment fields are copied from
\noderef{TashkinovHierarchy}.  This checks the new hierarchy.

For completeness, check all five fields of
\noderef{GoodHierarchy} with the same Gamma pairs.  Pair cardinality and
pairwise disjointness are literal.  The stored sets $D_{n,j}$ are the
same because all stage missing-color sets are unchanged.  Avoidance is
unchanged because every segment edge identity through the new terminal
prefix has the same color by (9)--(10).  Unique refresh is unchanged
because both the ordered stage differences and their missing sets are
unchanged.  The parallel-base field is unchanged because the base core,
closure, their common missing colors, and the relevant colored boundary
identities are fixed by the two constituent stability relations.  Each
minus-closed field concerns an old ordinary stage and is unchanged for
the same reason.  Equivalently these are precisely the stagewise
equalities encoded by \noderef{SameGoodHierarchy}.  Thus the new
hierarchy is genuinely good with the same Gamma pairs.                  \tag{12}

Equations (4), (11), and (12) give the actual non-elementarity, ETT,
maximum, base, hierarchy, and good-hierarchy fields of a counterexample
$C'$.  Its divider count is exactly $q$.  Equation (5) contradicts
minimality.  Hence $K_a=K_b$ in Location 2.  Alternative 3 never occurs
here because it requires $\alpha$ to be missing on $A$.

\textbf{Location 3: $a\in V(A)$ and $b\notin V(A)$.}
The conditional assumptions again give (a)--(b).  The first-missing
definition and Tashkinov augmentation rule give
\[
 \alpha\notin\varphi_n\langle T(b)-A\rangle,
 \qquad
 \beta\notin\overline\varphi_n(T(b)-b),                       \tag{13}
\]
and no displayed edge of $T(b)-A$ has color $\beta$, except possibly
the exact first connecting edge when $q=0$, $A=T_n$, and
$\beta\in D_n$.  This is obtained step by step from
\noderef{TashkinovAugmentation}: an edge added before $b$ is colored by
a color missing on its preceding prefix, whereas $b$ is the first
vertex missing $\beta$.

Suppose first that $\alpha$ is closed on $A$.  If $q>0$, the
last-divider instance of \noderef{step1} makes $\alpha,\beta$
$A$-interchangeable.  If $q=0$, use conclusion (iii) of
\noderef{interchange} in PE, or the terminal-ladder conclusion of
\noderef{TechnicalTheoremBase} and the lower-rung induction interface in
the non-PE endpoint.  Thus $K_a$ is the only nontrivial path component
meeting $A$.  If $K_a\ne K_b$, then $K_b$ is vertex-disjoint from $A$
by \noderef{TwoColorComponentsDisjoint}.  Swap $K_b$.  The outside-chain
lemma \noderef{OutsideKempeRecoloringStable}, (13), and the exact
construction (10)--(12) produce a same-divider counterexample on
$T(b)$, contradicting (5).  Therefore $K_a=K_b$.  Moreover the cited
interchangeability is exactly the additional conclusion required in
alternative 3.                                                        \tag{14}

We may henceforth assume
\[
        \alpha\text{ is not closed on }A,\qquad K_a\ne K_b,  \tag{15}
\]
and derive a contradiction under alternatives 1 and 2.

First let $q=0$.  Alternative 2 implies $\alpha\notin D_n$: this is
immediate from pair avoidance, while if $\alpha$ is missing on $T_n$ it
is removed from $D_n$ by \noderef{TashkinovStoredColors}.  If the last
extension were SE or RE, its constructor in
\noderef{TashkinovSeriesStep} would make every missing color outside
$D_n$ closed on $T_n=A$ (in SE use the selected stability and
\noderef{sc2}); this contradicts (15).  Hence the last extension is PE
and \noderef{HierarchyBase} gives the exact ordered union
\[
                         A=T_n\vee R_n.                       \tag{16}
\]

We now supply the elementarity premise required by the PE closure
argument.  If $n>1$, restrict the given ETT and maximum property through
rung $n-1$ and apply lower assertion (i).  If $n=1$, apply the
elementarity conclusion of \noderef{TechnicalTheoremBase}.  In either
case the pre-step coloring is elementary on $T_n$.  Thus
\noderef{ParallelCommonMissingClosed} applies with its genuine
pre-step-elementarity hypothesis.  If $a$ belonged to
$V(T_n)\cap V(R_n)$, then $\alpha$ would be common-missing on the two
constituents and the cited lemma would close it on their exact union,
contrary to (15).  Consequently
\[
 a\in V(T_n)\setminus V(R_n)\quad\hbox{or}\quad
 a\in V(R_n)\setminus V(T_n).                                \tag{17}
\]
This is where pre-step elementarity is used; PE alone has not been used
to infer closure.

In the first alternative of (17), conclusion (ii) of
\noderef{interchange} makes $K_a$ the only path component meeting
$T_n$; in the second, conclusion (i) does the same for $R_n$.  Thus the
different component $K_b$ is disjoint from the corresponding
constituent.  On the other constituent, (1), the fact that $a$ is its
unique possible $\alpha$-missing vertex in the core, and the fact that
$b$ is the first $\beta$-missing vertex outside the core show that every
vertex of $K_b$ in that constituent has both colors present.  A full
component swap therefore preserves its missing-color set.  Its incident
edge clause is preserved because any changed incident edge has color
$\alpha$ or $\beta$, neither of which is controlled there.  Hence the
swap is stable separately on $T_n$ with control $D_n$ and on $R_n$ with
empty control: it satisfies \noderef{CombinedStableColoring}.           \tag{18}

Equation (13) removes the possible first-edge exception: if it occurred,
the edge would have color $\beta\in D_n$, whereas alternative 2 and
$a\in A$ force either pair avoidance or $\alpha$ already missing on
$T_n$; in both cases the relevant conclusion of
\noderef{interchange} puts $K_b$ away from that exact connecting edge.
Apply \noderef{CombinedStableClosurePreserved} to (18).  It fixes every
displayed edge identity of both $T_n$ and $R_n$, and it makes the same
ordered tree $R_n$ a $\sigma_n$-closure of the same exact support prefix.
Thus every displayed identity of $T(b)$ retains its value.  Copy the
augmentation sequence edge by edge as in (11).  Apply \noderef{Td} to
the $T_n$ part of (18), obtain the corresponding ETT from
\noderef{ModColoring}, and transfer maximum property with
\noderef{MaximumPropertyTransport}.  The transported closure conclusion
and (16) supply the same base core.

For the zero-divider good hierarchy, pair size and disjointness are
literal, avoidance follows from the unchanged prefix edges, and the PE
base clause follows from the constituent missing-set and boundary-edge
equalities in (18); there are no refresh or positive-stage minus-closed
fields.  Thus the actual candidate satisfies every field, is
non-elementary by (4), and has the measure comparison (5).  This
contradicts minimality and settles $q=0$.                               \tag{19}

It remains to treat $q>0$.  Since $a\in A=T_{n,q}$, the color $\alpha$
is missing on $T_{n,q}$, so \noderef{rutcor} applies.  Its pre-step
elementarity premise is supplied exactly as in the paragraph preceding
(17), from lower assertion (i) or
\noderef{TechnicalTheoremBase}; hence this use also respects the PE
common-missing restriction.

Suppose first that some starred stage is $\alpha$-closed, and choose the
largest such index $r$.  By (15), $r<q$.  The first conclusion of
\noderef{rutcor} gives the exact edge-index exclusion
\[
 \alpha\notin
       \varphi_n\langle T_{n,q}-T_{n,r}^{*}\rangle.            \tag{20}
\]
If $r>0$, \noderef{step1} makes the two colors
$T_{n,r}^{*}$-interchangeable.  If $r=0$, use the exact PE/core
conclusion of \noderef{interchange} when the terminal step is PE, and
the lower-rung or \noderef{TechnicalTheoremBase} terminal-ladder
conclusion when it is not.  Therefore $K_b$ is disjoint from
$T_{n,r}^{*}$.

If $r>0$, truncate the current terminal augmentation sequence at the
exact tree $T_{n,r}^{*}=T_{n,r}$.  This is a same-rung ETT whose terminal
tree is $T_{n,r}$; it carries the first $r-1$ dividers, because its
terminal stage has index $(r-1)+1$.  Its ladder series, stored sets,
connecting-edge identities, and extension kinds are unchanged, so every
field of maximum property restricts.  Pair size, avoidance,
disjointness, unique refresh, the PE base field, and the minus-closed
fields through index $r-1$ are literal restrictions of the old good
hierarchy.  Since $K_b$ is disjoint from $T_{n,r}^{*}$, its swap is an
outside recoloring of this truncated ETT.  Apply
\noderef{HierarchyLEMStable} with the stated six lower-rung assertions.

If $r=0$, do not assert that the starred core is itself a same-rung ETT.
In the PE case $K_b$ is outside the exact core
$T_{n,0}^{*}=T_n\vee R_n$, so
\noderef{OutsideKempeRecoloringStable} restricts separately to $T_n$ and
$R_n$ and gives \noderef{CombinedStableColoring}; then
\noderef{CombinedStableClosurePreserved} transports the exact closure.
In the non-PE case the core is $T_n$, and the same outside-recoloring
lemma gives terminal-ladder stability directly.  In either case apply
\noderef{Td}; it uses the stated six lower-rung assertions and also
covers the one-rung endpoint.  Its universal mod-coloring clause supplies
the corresponding ETT data, and
\noderef{MaximumPropertyTransport} supplies maximum property.  In all
three subcases we have
a corresponding certified prefix with maximum property, the same base,
the same retained stages and Gamma pairs, and every retained
minus-closed field.                                                     \tag{21}

Combine (20) with (13).  Their adjacent exact index intervals show that
$\alpha$ is unused on $T(b)-T_{n,r}^{*}$.  Before $b$, no new edge has
color $\beta$, by the augmentation argument following (13).  Thus the
same exact augmentation steps extend the transported ETT in (21) to
$T(b)$.  Use the universal mod-coloring conclusion supplied through
\noderef{Td} to obtain the corresponding full ETT and
\noderef{MaximumPropertyTransport} for maximum property.

We check the good hierarchy after $r$, where (21) alone does not suffice.
The Gamma pairs themselves, their sizes, and disjointness are unchanged.
The stage missing sets are unchanged because $b$ is later than every
divider and every earlier vertex of $K_b$ has both colors present.
Avoidance and unique refresh are unchanged because the displayed
segment edges and stage-difference missing sets are unchanged.  For a
stage $j>r$, maximality of $r$ says that $\alpha$ is not closed there.
We also need, and now prove, that it is missing there.  Start at $j=q$,
where $a\in T_{n,q}$.  Descend one stage at a time.  If $\alpha$ were
missing on $T_{n,j+1}$ but not on $T_{n,j}^{*}$, its first witness would
lie in the vertex difference.  Elementarity from (1) and
\noderef{GoodHierarchyNewMissingClosed} would make it closed on
$T_{n,j+1}^{*}=T_{n,j+1}$, contradicting maximality of $r$.  Hence it
was already missing on $T_{n,j}^{*}$.  This descending induction proves
missingness at every index above $r$.  The minus-closed field of
\noderef{GoodHierarchy} would close any missing color outside its
reserved union.  Since $\alpha$ is not closed, it follows that
$\alpha$ lies in
$\bigcup_{\eta\in D_{n,j}}\Gamma^{j-1}_\eta$; $\beta$ is not missing on
that stage because its first missing vertex is $b$.  Hence changing an
$\alpha$- or $\beta$-colored boundary identity cannot violate the
minus-closed assertion: $\alpha$ is among the excluded reserved colors
and $\beta$ is not among the colors to which minus-closedness applies.
All other boundary colors are fixed.  This verifies the fifth field
stage by stage with the same Gamma pairs.  The parallel-base field lies
at or before $r$ and was transported in (21).                              \tag{22}

The data (21)--(22), (4), and the copied terminal sequence constitute a
same-divider hierarchy counterexample on $T(b)$.  Its actual measure is
smaller by (5), contradicting minimality.

Finally suppose no starred stage is $\alpha$-closed.  Use every part of
the second conclusion of \noderef{rutcor}.  For every $1\leq j\leq q$,
\[
 \alpha\in\bigcup_{\eta\in D_{n,j}}\Gamma^{j-1}_\eta
       \subseteq\Gamma^{j-1};                                \tag{23}
\]
the last extension is PE; $a\in V(T_n)\setminus V(R_n)$; and
\[
 \alpha\notin
       \varphi_n\langle T_{n,q}-T_n\rangle.                   \tag{24}
\]
By (1), no vertex of $R_n$ misses $\alpha$, since such a vertex would be
different from $a$ in the elementary core.  No vertex of $R_n$ misses
$\beta$, because $b$ is its first missing vertex and lies outside $A$.
Conclusion (ii) of \noderef{interchange} makes $K_b$ disjoint from
$T_n$.  Consequently the component swap is stable on $T_n$ and,
separately, on $R_n$: on $T_n$ it changes no incident identity, and on
$R_n$ every encountered component vertex has both colors present, so
the swap preserves its missing set.  This gives the exact
\noderef{CombinedStableColoring} premise, not an unsupported assertion
that PE alone supplies closure.                                        \tag{25}

Equations (24), (13), and the PE-base edge interval cover
$T(b)-T_n$.  The closure sequence of $R_n$ and elementarity of the core
exclude $\alpha$ from each exact edge of $R_n-T_n$; otherwise the
augmentation supplying that edge would name a second core vertex missing
$\alpha$.  The same first-missing argument excludes $\beta$.  Apply
\noderef{CombinedStableClosurePreserved} to (25).  It fixes every
displayed identity of $T_n$ and $R_n$ and transports the exact closure.
Together with the stated exclusions, all displayed edge identities of
$T(b)$ are fixed.  Apply \noderef{Td} to
the first stability relation in (25), copy the exact PE closure and
terminal augmentation sequences, obtain the corresponding ETT from
\noderef{ModColoring}, and transport maximum property by
\noderef{MaximumPropertyTransport}.                                   \tag{26}

Copy the hierarchy and Gamma pairs.  Pair size, disjointness, avoidance,
unique refresh, and the PE base field are preserved by (25)--(26).
For every positive stage, (23) places $\alpha$ in the exact reserved
union excluded from its minus-closed requirement, while $\beta$ is not
missing before $b$; all other boundary colors are unchanged.  Hence the
same-Gamma minus-closed field holds stage by stage.  We have again
constructed the actual base, hierarchy, good witness, corresponding
maximum ETT, and non-elementarity witness required by
\noderef{HierarchyCounterexample}.  The smaller measure (5) contradicts
minimality.                                                            \tag{27}

All possibilities under (15) contradict minimality, so $K_a=K_b$ in
Location 3.  Together with Locations 1 and 2, this proves the component
equality in alternatives 1 and 2.  Under alternative 3, Location 2 is
impossible, and (7), (8), or (14) gives
\noderef{TreeInterchangeable} on $A$.  By its definition, every two
nontrivial path components meeting $A$ are equal; since the common
component through $v_\alpha,v_\beta$ is one of them, it is the unique
one.  This proves alternative 3 and completes the proof.
\end{proof}
